\section*{Capítulo \ref{ch:b3:conservation-biology} --- Biologia da conservação}

\begin{solution}{exo:b3:conservation-biology:1}
\omterm{def:b3:conservation-biology:small}{Estocasticidade demográfica}: o acaso nos nascimentos, nas mortes e nos
sexos de uns poucos indivíduos --- o kakapo com cinquenta e um, com mais
machos que fêmeas. \omterm{def:b3:conservation-biology:small}{Estocasticidade ambiental}: um ano ruim que atinge
todos os indivíduos de uma vez --- uma seca, uma doença, um inverno
rigoroso golpeando uma população de algumas centenas. \omterm{def:b3:conservation-biology:small}{Efeito Allee}: o
crescimento per capita caindo em densidade baixa --- o pombo-passageiro,
cujas colônias não conseguiam se reproduzir depois de rareadas, e cujos
predadores restantes já não eram saturados.
\end{solution}

\begin{solution}{exo:b3:conservation-biology:2}
O tamanho de uma população ideal --- constante, com acasalamento
aleatório, sexos iguais e tamanhos de família de Poisson --- que perderia
heterozigosidade na mesma velocidade que a real. Números desiguais de
machos e fêmeas reprodutores; flutuações de tamanho, que pesam as
gerações menores; variância no tamanho da família acima da de Poisson,
com poucos indivíduos produzindo a maior parte dos filhotes (e também
acasalamento não aleatório e gerações sobrepostas).
\end{solution}

\begin{solution}{exo:b3:conservation-biology:3}
$S = cA^{z}$. Fração mantida $= (A'/A)^{z}$: $0.5^{0.25} = 0.84$;
$0.1^{0.25} = 0.56$; $0.01^{0.25} = 0.32$ --- um centésimo da área ainda
mantém, no fim, um terço das espécies.
\end{solution}

\begin{solution}{exo:b3:conservation-biology:4}
$N_{e} \ge 50$ mantém a subida da endogamia abaixo de $1/100$ por
geração, a taxa em que os criadores de animais acham a depressão
tolerável: protege contra a perda de aptidão no curto prazo. $N_{e} \ge
500$ equilibra a perda de variação pela deriva contra seu fornecimento
pela mutação, de modo que a população mantém a variação quantitativa de
que precisa para se adaptar: protege o longo prazo. As populações reais
precisam de cinco a dez vezes esses números no censo.
\end{solution}

\begin{solution}{exo:b3:conservation-biology:5}
$N_{e} = 4\times 8\times 120/128 = 30$; a heterozigosidade cai $1/60 =
\qty{1.7}{\%}$ por geração; um quarto se perde quando $(1 - 1/60)^{t} =
0.75$, $t = \ln 0.75/\ln(59/60) = 17$ gerações.
\end{solution}

\begin{solution}{exo:b3:conservation-biology:6}
$1/N_{e} = \tfrac{1}{5}(4/1000 + 1/20) = 0.0108$, $N_{e} = 93$: uma
geração com vinte fez cinco gerações se comportarem como noventa e três.
Heterozigosidade retida $(1 - 1/186)^{5} = 0.973$, contra $(1 -
1/2000)^{5} = 0.9975$ para um milhar constante.
\end{solution}

\begin{solution}{exo:b3:conservation-biology:7}
$F = 1 - (1 - 1/20)^{5} = 0.23$; aptidão reduzida em $5\times 2.3 = \qty{11}{\%}$. Todo filhote de uma fêmea texana com um macho da Flórida
tem $F = 0$; se as oito recém-chegadas forem dois quintos das fêmeas
reprodutoras, dois quintos da geração seguinte são exogâmicos e o $F$
médio cai a cerca de $0.6\times 0.23 = 0.14$ em uma geração --- que é o
que a recuperação mostrou.
\end{solution}

\begin{solution}{exo:b3:conservation-biology:8}
$r = \ln(100/31)/8 = \qty{0.15}{yr^{-1}}$. Continuando até 2020:
$100\,\mathrm{e}^{0.15\times 17} \approx 1200$. O parque abrigou cerca de
cem: os territórios se encheram, os wapitis declinaram, as alcateias
mataram membros umas das outras, e a sarna e a cinomose se espalharam ---
a dependência da densidade se instalou quando a área vazia foi ocupada.
\end{solution}

\begin{solution}{exo:b3:conservation-biology:9}
(a) Um declínio de \qty{60}{\%} em dez anos ultrapassa \qty{50}{\%}: Em
Perigo. (b) $180 < 250$ indivíduos maduros: Criticamente em Perigo. (c)
Área de distribuição de \qty{3000}{km^2}, abaixo de \qty{5000}{}: Em
Perigo. (d) \num{30000} indivíduos e um declínio de \qty{10}{\%}: nenhum
limiar atingido --- não ameaçada (no máximo Quase Ameaçada, a observar).
\end{solution}

\begin{solution}{exo:b3:conservation-biology:10}
Todos os $N$ casais falham com probabilidade $(1/2)^{N}$: $0.25$,
$0.031$, $0.001$, $10^{-6}$. O risco demográfico cai exponencialmente com
$N$ e é desprezível além de algumas dezenas de indivíduos; acima disso, o
risco que resta é ambiental, que não cai tão depressa, porque atinge
todos os indivíduos juntos.
\end{solution}

\begin{solution}{exo:b3:conservation-biology:11}
Um fragmento de um décimo: $0.1^{0.25} = 0.56$ das espécies originais.
Dez fragmentos abrigam entre $0.56$ (se todos abrigarem as mesmas
espécies) e, em princípio, mais que o todo (se cada um abrigasse espécies
diferentes --- impossível além do conjunto regional). Onde eles caem
depende da rotatividade entre fragmentos: hábitats diferentes em
fragmentos diferentes empurram o total para cima; as espécies que
precisam de grandes áreas, de hábitat de interior, ou que são perdidas de
todo fragmento pelos mesmos \omterm{prop:b3:conservation-biology:area}{efeitos de borda}, o empurram para baixo; e a
dívida de extinção de cada fragmento é paga separadamente, de modo que o
total de longo prazo costuma ficar abaixo do total do todo.
\end{solution}

\begin{solution}{exo:b3:conservation-biology:12}
A deriva retira heterozigosidade à taxa $1/2N_{e}$ por geração; os
imigrantes substituem uma fração $m$ do patrimônio genético por alelos
vindos de outro lugar. Com $mN_{e} \approx 1$, $m \approx 1/N_{e}$, o
dobro da taxa de perda, de modo que a variação é reposta mais depressa do
que decai e a população fica perto da fonte em frequência alélica
($F_{ST} = 1/(1
+ 4N_{e}m) \approx 0.2$). O custo: os alelos favorecidos
localmente são diluídos em $m$ por geração, de modo que a adaptação local
só sobrevive para características cujo coeficiente de seleção ultrapasse
cerca de $1/N_{e}$; as diferenças locais fracamente selecionadas são
engolidas.
\end{solution}

\begin{solution}{pb:b3:conservation-biology:1}
\textbf{1.} $N_{e} = 4\times 20\times 40/60 = 53$, contra um censo de 60.
\textbf{2.} $1/N_{e} = \tfrac{1}{5}(1/500 + 1/100 + 1/60 + 1/51 + 1/60) =
0.0130$, $N_{e} = 77$; as gerações de gargalo, de 51 e 60, dominam, e a
de 500 quase não conta.
\textbf{3.} $F = 1 - (1 - 1/154)^{5} = 0.032$; linear, $5/154 = 0.032$.
\textbf{4.} $6\times 0.32 = \qty{1.9}{\%}$.
\textbf{5.} $(1 - 1/154)^{5} = 0.968$: \qty{96.8}{\%} retida.
\textbf{6.} Dez gerações com $N_{e} = 53$: incremento $1 - (1 -
1/107)^{10} = 0.090$, total $F = 1 - 0.968\times 0.910 = 0.12$; perda de
aptidão $6\times 1.2 = \qty{7}{\%}$; um século.
\textbf{7.} Os casais nativo--nativo são $(60/70)^{2} = 0.73$ dos
acasalamentos e seus filhotes mantêm $F \approx 0.03$ (mais um pequeno
incremento); todos os demais têm $F = 0$: em média, $F \approx 0.73\times 0.035 = 0.026$, uma queda de um quarto em uma geração.
\textbf{8.} $F$ é a identidade por descendência dentro de um indivíduo, e
qualquer filhote exogâmico o zera; a heterozigosidade é uma propriedade
das frequências alélicas, e os alelos dos migrantes só se espalham à
medida que são transmitidos. Numa população ainda pequena, a deriva os
retirará de novo em dezenas de gerações, de modo que o resgate precisa
ser repetido ou a população precisa crescer.
\textbf{9.} $N_{e} = N$ com sexos iguais: 500 adultos. De 70 a $r =
0.05$: $\ln(500/70)/0.05 = 39$ anos.
\textbf{10.} A população ideal tem tamanhos de família de Poisson, com
variância igual à média de dois; $N_{e} = 4N/(2 + V_{k})$, de modo que
fazer cada casal contribuir igualmente ($V_{k} = 0$) dá $N_{e} = 2N$:
nenhuma linhagem se perde pela sorte do ninho.
\textbf{11.} Colete e armazene seu sêmen para inseminação artificial e o
use em várias fêmeas não aparentadas ao longo dos anos --- o custo é que
seus alelos poderiam passar a dominar um conjunto pequeno, de modo que
sua participação precisa ser limitada perto de $1/N_{e}$; ou
criopreserve tecido para clonagem ou produção de gametas mais tarde, ao
custo do atraso e do risco técnico. Não fazer nada perde os alelos da
linhagem para sempre.
\textbf{12.} \qty{160}{} milhões de dólares por umas 430 aves: cerca de
\num{370000} por ave. Caro por papagaio, barato contra uma espécie --- e
menos que uns poucos quilômetros de rodovia, que é a comparação que os
orçamentos de fato fazem.
\textbf{13.} $(300/2000)^{0.25} = 0.62$: cerca de 112 espécies das 180
mantidas, uma dívida de extinção de 68.
\textbf{14.} Cada fragmento, $(100/2000)^{0.25}\times 180 = 85$ espécies.
Inteiramente diferentes: até 255, impossível acima do conjunto de 180;
idênticas: 85. O que decide é quão diferentes são os hábitats dos
fragmentos e quantas espécies conseguem persistir em \qty{100}{km^2};
realisticamente, entre 85 e 112.
\textbf{15.} $(320/2000)^{0.25}\times 180 = 114$, contra 85 para três
fragmentos isolados idênticos: o corredor permite que a recolonização e o
fluxo gênico façam três populações pequenas se comportarem como uma
maior.
\textbf{16.} $N_{e} = 50$ com sexos iguais são 25 casais:
\qty{1250}{km^2}. Os \qty{300}{km^2} inteiros abrigam seis casais, $N_{e} \approx 12$; um fragmento, dois. Nenhum basta: o predador se perde
qualquer que seja a contagem de espécies, e proteger a área de que ele
precisa protege tudo o que é menor --- uma espécie guarda-chuva.
\textbf{17.} $\ln(200/10)/0.03 = 100$ anos: a dívida é paga ao longo de um
século, muito depois de a floresta ter sido cortada.
\textbf{18.} As bordas retiram hábitat de interior de todo fragmento, de
modo que a área utilizável é menor que a mapeada e a perda é maior que a
lei prevê; uma matriz de vegetação secundária ou de cercas vivas que
algumas espécies conseguem usar acrescenta área efetiva e torna a perda
menor.
\textbf{19.} $r = \ln(100/31)/8 = \qty{0.146}{yr^{-1}}$; tempo de
duplicação $\ln 2/0.146 = 4.7$ anos.
\textbf{20.} $N(8) = 120/[1 + (120/31 - 1)\mathrm{e}^{-1.17}] = 120/(1 +
2.87\times 0.31) = 63$. Os 100 observados são o próprio valor da
exponencial --- $r$ foi ajustado a esses mesmos dois pontos --- contra
os 63 da logística: nos primeiros anos os lobos não encontraram
dependência de densidade --- territórios vazios e wapitis abundantes.
\textbf{21.} $\ln(4000/\num{17000})/20 = -0.072$: \qty{7}{\%} ao ano. Cem
lobos tirando \num{2000} wapitis por ano de um rebanho de \num{8000} em
média são um quarto por ano --- mais que suficiente para explicar o
declínio sozinhos, embora a caça fora do parque, a seca e a predação de
filhotes por ursos-pardos o tenham compartilhado.
\textbf{22.} Sessenta indivíduos maduros, abaixo de 250: Criticamente em
Perigo. O predador com $N_{e} \approx 12$ (algumas dezenas de animais):
Criticamente em Perigo também.
\textbf{23.} De $3\times 10^{9}$ a $1$ em 44 anos: $r = -\ln(3\times
10^{9})/44 = -0.50$ por ano, caindo à metade a cada dezessete meses.
Nenhuma caça remove metade de bilhões todo ano; a queda foi mais lenta no
início e muito mais rápida no fim, porque um reprodutor colonial abaixo
de uma densidade-limiar para de se reproduzir enquanto seus predadores já
não são saturados --- o \omterm{def:b3:conservation-biology:small}{efeito Allee} transformou um declínio num colapso.
\textbf{24.} Em 1850 a população estava na casa dos bilhões, sua taxa de
crescimento medida era positiva e sua variância, pequena, e uma análise
de viabilidade projeta a dinâmica que lhe é dada: ela teria posto o risco
em zero. Ela teria deixado passar o forçamento externo --- as ferrovias e
o telégrafo, que deixaram os caçadores seguir e esvaziar todo dormitório,
e o desmatamento --- e o limiar de Allee num reprodutor colonial, nenhum
dos quais está na demografia de uma espécie abundante; uma AVP extrapola,
e as causas da extinção costumam ser novas.
\textbf{25.} $N_{e} \approx 53$ hoje; $F \approx 0.12$ após mais dez
gerações com 60 (\qty{7}{\%} da aptidão); a floresta encolhida mantém
cerca de \qty{62}{\%} de suas espécies.
\end{solution}
